\documentclass[10pt,a4paper]{amsart}
\usepackage{amsmath,amssymb,amsthm,graphicx}
% For commutative diagrams we use Xy-pic
\usepackage[all]{xy}

% We use 2cell for 2-commutative diagrams.
\xyoption{2cell}
% Theorem environments.
%
\theoremstyle{plain}
\newtheorem{theorem}[subsection]{Theorem}
\newtheorem{proposition}[subsection]{Proposition}
\newtheorem{lemma}[subsection]{Lemma}

\theoremstyle{definition}
\newtheorem{definition}[subsection]{Definition}
\newtheorem{example}[subsection]{Example}
\newtheorem{exercise}[subsection]{Exercise}
\newtheorem{situation}[subsection]{Situation}

\theoremstyle{remark}
\newtheorem{remark}[subsection]{Remark}
\newtheorem{remarks}[subsection]{Remarks}

\numberwithin{equation}{subsection}

% Macros
%
\def\lim{\mathop{\mathrm{lim}}\nolimits}
\def\colim{\mathop{\mathrm{colim}}\nolimits}
\def\Spec{\mathop{\mathrm{Spec}}}
\def\Hom{\mathop{\mathrm{Hom}}\nolimits}
\def\Ext{\mathop{\mathrm{Ext}}\nolimits}
\def\SheafHom{\mathop{\mathcal{H}\!\mathit{om}}\nolimits}
\def\SheafExt{\mathop{\mathcal{E}\!\mathit{xt}}\nolimits}
\def\Sch{\mathit{Sch}}
\def\Mor{\mathop{\mathrm{Mor}}\nolimits}
\def\Ob{\mathop{\mathrm{Ob}}\nolimits}
\def\Sh{\mathop{\mathit{Sh}}\nolimits}
\def\NL{\mathop{N\!L}\nolimits}
\def\CH{\mathop{\mathrm{CH}}\nolimits}
\def\proetale{{pro\text{-}\acute{e}tale}}
\def\etale{{\acute{e}tale}}
\def\QCoh{\mathit{QCoh}}
\def\Ker{\mathop{\mathrm{Ker}}}
\def\Im{\mathop{\mathrm{Im}}}
\def\Coker{\mathop{\mathrm{Coker}}}
\def\Coim{\mathop{\mathrm{Coim}}}

% Boxtimes
%
\DeclareMathSymbol{\boxtimes}{\mathbin}{AMSa}{"02}

%
% Macros for moduli stacks/spaces
%
\def\QCohstack{\mathcal{QC}\!\mathit{oh}}
\def\Cohstack{\mathcal{C}\!\mathit{oh}}
\def\Spacesstack{\mathcal{S}\!\mathit{paces}}
\def\Quotfunctor{\mathrm{Quot}}
\def\Hilbfunctor{\mathrm{Hilb}}
\def\Curvesstack{\mathcal{C}\!\mathit{urves}}
\def\Polarizedstack{\mathcal{P}\!\mathit{olarized}}
\def\Complexesstack{\mathcal{C}\!\mathit{omplexes}}
% \Pic is the operator that assigns to X its picard group, usage \Pic(X)
% \Picardstack_{X/B} denotes the Picard stack of X over B
% \Picardfunctor_{X/B} denotes the Picard functor of X over B
\def\Pic{\mathop{\mathrm{Pic}}\nolimits}
\def\Picardstack{\mathcal{P}\!\mathit{ic}}
\def\Picardfunctor{\mathrm{Pic}}
\def\Deformationcategory{\mathcal{D}\!\mathit{ef}}
\usepackage{cleveref}
\crefname{originallemma}{original lemma}{original lemmas}
\Crefname{originallemma}{Original lemma}{Original lemmas}
\makeatletter
\newlabel{divisors-definition-embedded}{{05AK}{source}{}{sourcepointer-1}{}}
\newlabel{divisors-definition-embedded@cref}{{[originallemma][0][]05AK}{[1][0][]source}{}{}{}}
\makeatother
\crefname{originallemma}{original lemma}{original lemmas}
\Crefname{originallemma}{Original lemma}{Original lemmas}
\makeatletter
\newlabel{divisors-lemma-components-locally-finite}{{0BE1}{source}{}{sourcepointer-2}{}}
\newlabel{divisors-lemma-components-locally-finite@cref}{{[originallemma][0][]0BE1}{[1][0][]source}{}{}{}}
\makeatother
\crefname{originallemma}{original lemma}{original lemmas}
\Crefname{originallemma}{Original lemma}{Original lemmas}
\makeatletter
\newlabel{lemma-additivity-divisors-restricted}{{02QG}{source}{}{sourcepointer-3}{}}
\newlabel{lemma-additivity-divisors-restricted@cref}{{[originallemma][0][]02QG}{[1][0][]source}{}{}{}}
\makeatother
\crefname{originallemma}{original lemma}{original lemmas}
\Crefname{originallemma}{Original lemma}{Original lemmas}
\makeatletter
\newlabel{lemma-rational-equivalence-family}{{02S4}{source}{}{sourcepointer-4}{}}
\newlabel{lemma-rational-equivalence-family@cref}{{[originallemma][0][]02S4}{[1][0][]source}{}{}{}}
\makeatother
\crefname{originallemma}{original lemma}{original lemmas}
\Crefname{originallemma}{Original lemma}{Original lemmas}
\makeatletter
\newlabel{situation-setup}{{02QL}{source}{}{sourcepointer-5}{}}
\newlabel{situation-setup@cref}{{[originallemma][0][]02QL}{[1][0][]source}{}{}{}}
\makeatother
\title{The two fibre cycle classes of a closed subscheme of a product with the projective line are rationally equivalent}
\author{The Stacks Project Authors}
\date{stacks-project:chow.tex:local-snapshot-20261012; GNU FDL 1.2 or later}
\begin{document}
\maketitle
\noindent\textit{Editorial note.} All delivered mathematical text is the author's own, copied from the pinned Stacks Project chapter file chow.tex; no statement or proof step was rewritten, reordered or removed.\par\medskip
\noindent\textit{Editorial note.} Preamble provenance: the pinned chapter file opens with \textbackslash{}input\{preamble\} and carries no preamble text of its own, so the project's single Stacks-wide preamble staged at sources/stacks-chow/preamble.tex is excerpted verbatim as two spans, lines 14-18 (Xy-pic setup) and lines 154-217 (the native theorem declarations and mathematical macros). Nothing was added, removed or redefined; the wrapper class and the amsmath/amssymb/amsthm/graphicx base packages are supplied by the assembler.\par\medskip
\noindent\textit{Editorial note.} Same-chapter pointers are located at their \textbackslash{}label lines inside the pinned chapter file; cross-chapter pointers are located in the project's own tag list staged in the snapshot (sources/stacks-chow/tags.txt), at the verbatim lines carrying their canonical tags.\par\medskip
\noindent\textit{Allowed standard prerequisites (editorial).}\begin{itemize}
\item Embedded points of a closed subscheme and the local finiteness of its components. Divisors, labels divisors-definition-embedded (statement, proof line 3413) and divisors-lemma-components-locally-finite (proof line 3429). Fixes the hypothesis on embedded points used at the generic point of a component of the fibre, and makes the family of irreducible components W\_i of Z of dimension k + 1 locally finite so that [Z]\_\{k+1\} is a well defined sum.
\item Rational equivalence of a family of fibres. Same chapter chow.tex, label lemma-rational-equivalence-family, invoked at proof line 3435. Turns the two coefficientwise identities into the asserted rational equivalence of [Z\_0]\_k and [Z\_infty]\_k.
\item Additivity of divisors for a restricted family. Same chapter chow.tex, label lemma-additivity-divisors-restricted, invoked at proof line 3466. Supplies the final identity comparing length\_A(A/(t,I)) with the sum over the minimal primes over I of the residue length times the local length of A/I.
\end{itemize}\medskip
\ifdefined\hypertarget\hypertarget{sourcepointer-1}{}\fi
\noindent\textit{Source reference.} Original lemma 05AK: sources/stacks-chow/tags.txt, source lines 6853--6853. Navigation pointer for \textbackslash{}ref\{divisors-definition-embedded\} used inside the counted span (line 3413); the argument that consumes it is retained in the span.\par\medskip
\ifdefined\hypertarget\hypertarget{sourcepointer-2}{}\fi
\noindent\textit{Source reference.} Original lemma 0BE1: sources/stacks-chow/tags.txt, source lines 14334--14334. Navigation pointer for \textbackslash{}ref\{divisors-lemma-components-locally-finite\} used inside the counted span (line 3429); the argument that consumes it is retained in the span.\par\medskip
\ifdefined\hypertarget\hypertarget{sourcepointer-3}{}\fi
\noindent\textit{Source reference.} Original lemma 02QG: sources/stacks-chow/pinned-original-source.tex, source lines 373--373. Navigation pointer for \textbackslash{}ref\{lemma-additivity-divisors-restricted\} used inside the counted span (line 3466); the argument that consumes it is retained in the span.\par\medskip
\ifdefined\hypertarget\hypertarget{sourcepointer-4}{}\fi
\noindent\textit{Source reference.} Original lemma 02S4: sources/stacks-chow/pinned-original-source.tex, source lines 3335--3335. Navigation pointer for \textbackslash{}ref\{lemma-rational-equivalence-family\} used inside the counted span (line 3435); the argument that consumes it is retained in the span.\par\medskip
\ifdefined\hypertarget\hypertarget{sourcepointer-5}{}\fi
\noindent\textit{Source reference.} Original lemma 02QL: sources/stacks-chow/pinned-original-source.tex, source lines 1372--1372. Navigation pointer for \textbackslash{}ref\{situation-setup\} used inside the counted span (line 3405); the argument that consumes it is retained in the span.\par\medskip
\begin{lemma}
\label{lemma-closed-subscheme-cross-p1}
Let $(S, \delta)$ be as in Situation \ref{situation-setup}.
Let $X$ be a scheme locally of finite type over $S$.
Let $Z$ be a closed subscheme of $X \times \mathbf{P}^1$.
Assume
\begin{enumerate}
\item $\dim_\delta(Z) \leq k + 1$,
\item $\dim_\delta(Z_0) \leq k$, $\dim_\delta(Z_\infty) \leq k$, and
\item for any embedded point $\xi$ (Divisors, Definition
\ref{divisors-definition-embedded}) of $Z$ either
$\xi \not \in Z_0 \cup Z_\infty$ or $\delta(\xi) < k$.
\end{enumerate}
Then $[Z_0]_k \sim_{rat} [Z_\infty]_k$ as $k$-cycles on $X$.
\end{lemma}

\begin{proof}
Let $\{W_i\}_{i \in I}$ be the collection of irreducible
components of $Z$ which have $\delta$-dimension $k + 1$.
Write
$$
[Z]_{k + 1} = \sum n_i[W_i]
$$
with $n_i > 0$ as per definition. Note that $\{W_i\}$
is a locally finite collection of closed subsets of
$X \times_S \mathbf{P}^1_S$ by
Divisors, Lemma \ref{divisors-lemma-components-locally-finite}.
We claim that
$$
[Z_0]_k = \sum n_i[(W_i)_0]_k
$$
and similarly for $[Z_\infty]_k$. If we prove this then the lemma
follows from Lemma \ref{lemma-rational-equivalence-family}.

\medskip\noindent
Let $Z' \subset X$ be an integral closed subscheme of $\delta$-dimension $k$.
To prove the equality above it suffices to show that the coefficient $n$
of $[Z']$ in $[Z_0]_k$ is the same as the coefficient $m$ of
$[Z']$ in $\sum n_i[(W_i)_0]_k$. Let $\xi' \in Z'$ be the generic point.
Set $\xi = (\xi', 0) \in  X \times_S \mathbf{P}^1_S$.
Consider the local ring $A = \mathcal{O}_{X \times_S \mathbf{P}^1_S, \xi}$.
Let $I \subset A$ be the ideal cutting out $Z$, in other words so that
$A/I = \mathcal{O}_{Z, \xi}$. Let $t \in A$ be the element cutting
out $X \times_S D_0$ (i.e., the coordinate of $\mathbf{P}^1$ at zero
pulled back). By our choice of $\xi' \in Z'$ we have $\delta(\xi) = k$
and hence $\dim(A/I) = 1$. Since $\xi$ is not an embedded point by
assumption (3) we see that $A/I$ is Cohen-Macaulay. Since $\dim_\delta(Z_0)
= k$ we see that $\dim(A/(t, I)) = 0$ which implies that $t$
is a nonzerodivisor on $A/I$. Finally, the irreducible closed subschemes
$W_i$ passing through $\xi$ correspond to the minimal primes
$I \subset \mathfrak q_i$ over $I$. The multiplicities $n_i$ correspond
to the lengths $\text{length}_{A_{\mathfrak q_i}}(A/I)_{\mathfrak q_i}$.
Hence we see that
$$
n = \text{length}_A(A/(t, I))
$$
and
$$
m = \sum
\text{length}_A(A/(t, \mathfrak q_i))
\text{length}_{A_{\mathfrak q_i}}(A/I)_{\mathfrak q_i}
$$
Thus the result follows from
Lemma \ref{lemma-additivity-divisors-restricted}.
\end{proof}

\end{document}
